\begin{figure}[ht]
\centering
\begin{tikzpicture}[scale=1.0]
% A modular clock for Z/12Z, the residues placed as on a clock face,
% with the orbit 0,5,10,3,8,1,6,11,4,9,2,7 of repeated +5 traced.
\def\rad{2.6}
\draw[thick, draw=enggray] (0,0) circle (\rad);
\foreach \k in {0,...,11} {
  \pgfmathsetmacro{\ang}{90 - 30*\k}
  \node[circle, draw=engblack, fill=englime!20, inner sep=0pt,
        minimum size=6.5mm, font=\small] (m\k) at (\ang:\rad) {$\k$};
}
% Trace the orbit of repeated addition of 5 modulo 12.
\foreach \a/\b in {0/5, 5/10, 10/3, 3/8, 8/1, 1/6, 6/11, 11/4,
                   4/9, 9/2, 2/7, 7/0} {
  \draw[->, thick, draw=engblue, shorten >=2pt, shorten <=2pt]
    (m\a) -- (m\b);
}
\node[font=\scriptsize, anchor=north, enggray] at (0, -\rad-0.45)
  {$+5 \pmod{12}$ visits every residue: $\gcd(5, 12) = 1$};
\end{tikzpicture}
\caption{Arithmetic in $\Z/12\Z$ drawn on a clock face. Repeated
addition of $5$, traced by the arrows, visits all twelve residues
before returning to $0$. The orbit closes after exactly $12$ steps
because $\gcd(5, 12) = 1$, so $5$ generates the additive group. An
increment coprime to the modulus is a full-period generator; an
increment sharing a factor with the modulus traces a proper
sub-cycle and never reaches every residue. The same fact governs
the choice of stride in a linear-congruential pseudo-random
generator and the choice of probe increment in open-addressing
hash tables.}
\label{fig:vol02:ch17:modular-clock}
\end{figure}
